\documentclass{article}
\input{../../../lib.sty}

\title{Joint discrete-Gaussian group-by strategy for exact integer aggregates}
\date{}
\begin{document}
\maketitle

\section{Scope}

This argument concerns the optional \texttt{JointGaussian} strategy in
\texttt{joint\_gaussian.rs}. It does not replace the proof of
\texttt{make\_private\_group\_by}: the public constructor continues to use
independent composition. The strategy is eligible only when \emph{every}
materialized private aggregate is an implicitly scaled integer length or a
bounded, direct-column \texttt{Int32} sum. All other workloads fall back.

The existing scalar constructors enforce numeric validity and reject integer
sums that could overflow within the maximum group length. Consequently, each
eligible coordinate is an exact integral sum of per-record contributions
bounded by a public interval. The argument does not apply to floating-point
sum/variance, whose numerical relaxation requires separate analysis.

\section{Joint stability}

For each record \(x\), let \(v(x)\in\mathbb Z^d\) contain its contribution to
every released coordinate in one group. Length coordinates contribute \(1\).
A bounded sum coordinate contributes a clipped (and, for nulls, publicly
imputed) integer in its declared interval. Write
\[
 B_2=\sqrt{\sum_{j=1}^d \max(|a_j|,|b_j|)^2}.
\]
Thus \(\|v(x)\|_2\le B_2\) for every permitted row.

Suppose neighboring grouped datasets differ in at most \(l_0\) groups,
\(l_1\) rows total, and \(l_\infty\) rows in any single group. If group \(g\)
has \(m_g\) changed rows, then
\(\|\Delta q_g\|_2\le m_g B_2\) by the triangle inequality.
It follows that
\[
  \|\Delta q\|_2^2
  \le B_2^2\sum_g m_g^2
  \le B_2^2\min\bigl(
    l_1^2,\; l_\infty l_1,\; l_0 l_\infty^2
  \bigr).
\]
The implementation computes an upward-rounded bound on this quantity.
Every private leaf is included, and duplicates are removed by the preceding
workload planner. The bound is therefore on the \emph{full} private vector,
not an arbitrary subset of its coordinates.

\section{Gaussian release and accounting}

The function adds independent samples from the same discrete Gaussian
distribution of scale \(\sigma\) to each integer coordinate in each group.
The existing OpenDP scalar Gaussian mechanism certifies zCDP cost
\(\Delta_j^2/(2\sigma^2)\) for each coordinate with integer shift
\(\Delta_j\). Product composition over all coordinates and groups gives
\[
  \rho \le \frac{\|\Delta q\|_2^2}{2\sigma^2}.
\]
The implementation computes this expression using a conservative
floating-point distance bound, exact rational arithmetic for the quotient,
and upward rounding for the returned privacy loss.

For implicit calibration, \(\sigma\) is the upward-rounded product of the
requested global scale and the joint sensitivity at the smallest admissible
dataset distance (one added/removed row, or two for invariant lengths).
Every leaf receives the same computed scale via the existing scalar shim
compiler. The joint privacy map uses that same scale value.

For public-key joins, all missing-group coordinates are independently
imputed using the scalar constructors' noise fillers before deterministic
reconstruction. For private-key filters, the existing certified noisy-count
tail bound is applied separately as the catastrophic-event probability.
Any final arithmetic, filtering and projection are postprocessing.

When \(\sigma=0\) and the sensitivity is positive, the joint privacy map
returns infinite loss. When the input distance is zero, the map returns zero.
The strategy is only enabled for zCDP or approximate zCDP; other measures
retain the independent-composition strategy.

\end{document}
